\section[Estimates of $c_n$ via $V_n$]{Estimates of $\boldsymbol{c_n}$ via $\boldsymbol{V_n}$}\label{section4}

In all of this section we consider symmetric moment problems given by a positive sequence~$(b_n)$, and it is tacitly assumed that the moment problem is indeterminate in order to be able to define~$c_n$ from~\eqref{eq:c_k}, hence $\sum 1/b_n<\infty$ by a classical result due to Carleman, cf.~\cite[p.~24]{Ak}.

For $k,l\ge 0$ define
\begin{gather}\label{eq:vkl}
v_{k,l}:= (-1)^l b_0b_1\cdots b_{k-1} b_{k,k+2l},
\end{gather}
where $b_{k,k+2l}$ is the coefficient to $x^k$ in $P_{k+2l}$, so $(-1)^lb_{k,k+2l}>0$ and hence $v_{k,l}>0$.

Observe that $v_{k,0}=1$ and since $b_{k,j}=0$ if $j<k$ it makes sense to define $v_{k,l}=0$ if $l<0$.

From \eqref{eq:P2n(0)} we know
\begin{gather}\label{eq:v0l}
P_{2l}(0)=b_{0,2l}=(-1)^lv_{0,l}=(-1)^l\prod_{j=1}^l \frac{b_{2j-2}}{b_{2j-1}}.
\end{gather}

Notice that by \eqref{eq:c_k} and \eqref{eq:Vn}
\begin{gather}\label{eq:V_n1}
V_n=b_0b_1\cdots b_{n-1} c_n=\left(\sum_{l=0}^\infty v_{n,l}^2\right)^{1/2}.
\end{gather}

\begin{lem}\label{thm:vkl}
The coefficients $v_{k,l}$ satisfy
\begin{gather}\label{eq:vkl1} v_{k+1,l}={b_k\over b_{k+2l}}v_{k,l}+ {b_{k+2l-1}\over b_{k+2l}}v_{k+1,l-1},\qquad k,l\ge 0.
\end{gather}
\end{lem}

\begin{proof}We have
\begin{gather*}
x P_{k+2l}=b_{k+2l}P_{k+2l+1}+b_{k+2l-1}P_{k+2l-1}.
\end{gather*}
By equating the coefficients of $x^{k+1}$ we obtain
\begin{gather*}
b_{k,k+2l}=b_{k+2l}b_{k+1,k+2l+1}+ b_{k+2l-1}b_{k+1,k+2l-1}.
\end{gather*}
Multiplying the last identity by $(-1)^lb_0b_1\cdots b_{k-1}$ leads to
\begin{gather*}
 v_{k,l}={b_{k+2l}\over b_k}v_{k+1,l}- { b_{k+2l-1}\over b_k}v_{k+1,l-1}.
\end{gather*}
 This gives (\ref{eq:vkl1}).
\end{proof}

The following result follows easily from \eqref{eq:v0l} and \eqref{eq:vkl1}.

\begin{prop}\label{thm:compV} Assume $(b_n)$ and $\big(\tilde{b}_n\big)$ are positive sequences satisfying
\begin{gather*}
{b_{n-1}\over b_{n}} \le {\tilde{b}_{n-1}\over \tilde{b}_{n}} ,\qquad n\ge 1.
\end{gather*}
 Then $v_{k,l} \le \tilde{v}_{k,l} $ and $V_k \le \tilde{V}_k$.
\end{prop}

For $n_0\in\mathbb N$ the shifted sequence $b_n^{(n_0)}:=b_{n+n_0}$ also leads to an indeterminate moment problem, cf.~\cite[p.~28]{Ak}, and we shall consider the associated quantities $v_{k,l}^{(n_0)}$, $V_k^{(n_0)}$.

\begin{prop}\label{thm:shiftV1} The expressions above satisfy
\begin{gather}\label{eq:shiftV1}
v_{k,l}^{(1)}\le v_{k+1,l},\qquad V_k^{(1)}\le V_{k+1},\qquad k,l\ge 0,\\
\label{eq:shiftV2}
v_{k+1,l}\le L v_{k,l}^{(1)},\qquad V_{k+1}\le L V_k^{(1)},\qquad k,l\ge 0,
\end{gather}
where $L:=\sum\limits_{n=0}^\infty P_n^2(0)$. For $n_0\in\mathbb N$ we similarly have
\begin{gather}\label{eq:shiftV1a}
v_{k,l}^{(n_0)}\le v_{k+n_0,l}\le L(n_0) v_{k,l}^{(n_0)},\qquad V_k^{(n_0)}\le V_{k+n_0}\le L(n_0) V_k^{(n_0)},
\end{gather}
for a suitable constant $L(n_0)>0$. Furthermore,
\begin{gather} \label{eq:shiftV3}
\limsup_{k\to\infty} V_k^{1/k}=\limsup_{k\to\infty} \big(V_k^{(n_0)}\big)^{1/k}.
\end{gather}
\end{prop}

\begin{proof} From \eqref{eq:vkl1} we get
\begin{gather}\label{eq:shiftvkl1}
 v_{k+1,l}^{(1)}=\frac{b_{k+1}}{b_{k+2l+1}}v_{k,l}^{(1)}+ \frac{b_{k+2l}}{b_{k+2l+1}}v_{k+1,l-1}^{(1)},\qquad k,l\ge 0,\\
 \label{eq:vkl1+}
v_{k+2,l}=\frac{b_{k+1}}{b_{k+2l+1}}v_{k+1,l}+ \frac{b_{k+2l}}{b_{k+2l+1}}v_{k+2,l-1},\qquad k,l\ge 0.
\end{gather}

On the other hand
\begin{gather}\label{eq:v1l}
v_{1,l}=\frac{b_0}{b_{2l}} v_{0,l}+\frac{b_{2l-1}}{ b_{2l}}v_{1,l-1}\ge \frac{b_{2l-1}}{ b_{2l}}v_{1,l-1}.
\end{gather}
Iterating this we get using $v_{1,0}=1$
\begin{gather}\label{eq:v0l(1)}
v_{1,l}\ge \prod_{j=1}^l \frac{b_{2j-1}}{b_{2j}}=v_{0,l}^{(1)},
\end{gather}
where the last equality follows from \eqref{eq:v0l}. This proves the left inequality in~\eqref{eq:shiftV1} for $k=0$ and all $l$, and it is clear for $l=0$ and all~$k$. In particular it is true for $k+l\le 1$. Assume that it holds for $k+l\le n$ for some $n\ge 1$. To prove it for $k+l=n+1$, it suffices to prove $v_{k+1,l}^{(1)}\le v_{k+2,l}$ when $l\ge 1$, but this follows immediately by~\eqref{eq:shiftvkl1} and~\eqref{eq:vkl1+} and the induction hypothesis. The right inequality in~\eqref{eq:shiftV1} is an immediate consequence of the left inequality.

Defining $t_l=v_{1,l}/v_{0,l}^{(1)}$, $l\ge 0$ we get by \eqref{eq:v0l(1)}
\begin{gather*}
t_l=v_{1,l}\prod_{j=1}^l \frac{b_{2j}}{b_{2j-1}},\qquad l\ge 1,
\end{gather*}
so by \eqref{eq:v1l} and \eqref{eq:v0l}
\begin{gather*}
t_l=\left(\frac{b_0}{b_{2l}}v_{0,l}+\frac{b_{2l-1}}{b_{2l}}v_{1,l-1}\right)\prod_{j=1}^l \frac{b_{2j}}{b_{2j-1}}=\frac{b_0}{b_{2l}}\prod_{j=1}^l\frac{b_{2j-2}b_{2j}}{b_{2j-1}^2} + t_{l-1}.
\end{gather*}
We then get for $l\ge 1$
\begin{gather*}
t_l=t_0+\sum_{m=1}^l \frac{b_0}{b_{2m}}\prod_{j=1}^m \frac{b_{2j-2}b_{2j}}{b_{2j-1}^2} =1+\sum_{m=1}^l P_{2m}^2(0)\le \sum_{n=0}^\infty P_n^2(0)=L,
\end{gather*}
hence $v_{1,l}\le L v_{0,l}^{(1)}$ for all $l$. By induction as above we then get the left side of~\eqref{eq:shiftV2}, and the right side follows.

The inequalities \eqref{eq:shiftV1a} easily follow by iteration with a suitable constant $L(n_0)>0$ independent of~$k$, and \eqref{eq:shiftV3} follows.
\end{proof}

\begin{thm}\label{thm:V_nlogcv} Assume $(b_n)$ is $q$-increasing, cf.\ Definition~{\rm \ref{thm:defqinc}}. The quantity $V_n$ from \eqref{eq:Vn} satisfies
\begin{gather}\label{eq:V_nlogcv}
V_n\le\frac{1}{\big(q^2;q^2\big)_\infty\sqrt{1-q^2}}.
\end{gather}
This holds in particular if $(b_n)$ is log-convex and strictly increasing with $q=b_0/b_1<1$.
\end{thm}

\begin{proof} Note that the moment problem is indeterminate by Proposition~\ref{thm:qinc-indet}. From~\eqref{eq:vkl1} and \eqref{eq:logcv-q} we obtain
\begin{gather}\label{eq:vkl1q}
v_{k+1,l}\le q^{2l} v_{k,l}+qv_{k+1,l-1}.
\end{gather}

We claim that
\begin{gather}\label{eq:vklq}
v_{k,l}\le q^l\prod_{j=1}^l \big(1-q^{2j}\big)^{-1}=\frac{q^l}{\big(q^2;q^2\big)_l},\qquad k,l\ge 0,
\end{gather}
which is certainly true for $l=0$ since $v_{k,0}=1$ and for $k=0$ since by \eqref{eq:v0l}
\begin{gather*}
v_{0,l}\le q^l\le \frac{q^l}{\big(q^2;q^2\big)_l}.
\end{gather*}

 We shall proceed by induction in $k+l$, and as already observed \eqref{eq:vklq} holds when $k+l\le 1$.

Suppose \eqref{eq:vklq} holds for $k+l\le n$ for some $n\ge 1$ and we shall prove it for $k+l=n+1$. From \eqref{eq:vkl1q} and the induction hypothesis we get for $l\ge 1$
\begin{gather*}
v_{k+1,l} \le q^{2l}v_{k,l}+q\, v_{k+1,l-1} \le \frac{q^{3l}}{\big(q^2;q^2\big)_{l}}+\frac{q^{l}}{\big(q^2;q^2\big)_{l-1}}=\frac{q^{l}}{\big(q^2;q^2\big)_{l}}.
\end{gather*}

From \eqref{eq:vklq} we get
\begin{gather*}
v_{k,l}\le \frac{q^l}{\big(q^2;q^2\big)_\infty}
\end{gather*}
and finally
\begin{gather*}
V_n^2=\sum_{l=0}^\infty v_{n,l}^2\le \big(q^2;q^2\big)_\infty^{-2}\sum_{l=0}^\infty q^{2l},
\end{gather*}
leading to \eqref{eq:V_nlogcv}.
\end{proof}

The boundedness of $(V_n)$ also holds in a slightly more general case.

\begin{thm}\label{thm:Vevtlogconv} Assume $(b_n)$ is eventually $q$-increasing, cf.\ Definition~{\rm \ref{thm:defqinc}}. The sequence $(V_n)$ from \eqref{eq:Vn} is bounded.
\end{thm}

\begin{proof} Note that the moment problem is indeterminate by Proposition~\ref{thm:qinc-indet}. We assume that \eqref{eq:qinc} holds. Then the sequence $\big(b_n^{(n_0)}\big)$ is $q$-increasing, so $\big(V_n^{(n_0)}\big)$ is bounded by Theorem~\ref{thm:V_nlogcv}. The boundedness of $(V_n)$ follows from \eqref{eq:shiftV1a}.
\end{proof}

Combining Theorem~\ref{thm:evtlogconvex} and Theorem~\ref{thm:Vevtlogconv} we get:

\begin{cor}\label{thm:VnsqrtUn} Assume $(b_n)$ is eventually $q$-increasing. Then $c_n\sqrt{s_{2n}}=V_n\sqrt{U_n}$ is a bounded sequence.
\end{cor}
